\section{问题分析}

\subsection{针对问题一的分析}
问题一是在电价、小区负载与光伏出力均为已知固定日内规律的条件下，确定当天每 10\,min 的购电量与储能充放电策略。
由于各时段电价差异显著，而光伏出力高峰与负载高峰在时间上并不重合，核心思路是利用储能实现电能在不同时段之间的转移：
在低价与光伏富余时段充电，在电价较高时段放电以减少外网购电。
将一天离散为 144 个时段后，各时段须满足功率平衡，储能受储电量上下限、最大充放电功率、
充放电效率以及日循环（0:00 与 24:00 储电量相同）等约束；目标函数与全部约束均为决策变量的线性关系，
故该问题可转化为以全天购电费用最小为目标的线性规划问题，输出各时段购电量、储能充放电量与全天购电费用。

\subsection{针对问题二的分析}
问题二取消了“日内规律每天相同”的假设，负载与光伏逐日变化，且计划购电按照付不议结算，
供电不足部分需按当时电价的 5 倍紧急购电——购电过多会造成浪费，购电不足则承担高昂的紧急购电成本，
因此关键在于提高日前预测精度并合理权衡两者。
首先分析附件 2 的全年数据：小区负载呈现显著的星期型双状态特征（周日至周四高、周五周六低），
光伏则具有较强的季节性与相邻日持续性；据此采用“上周同星期实际值”作为预测基础\cite{fpp3}，
以避免持续性预测在星期状态切换日失效。
其次，由于紧急购电单价是正常电价的 5 倍，预测偏低的代价远高于偏高，
故依据报童模型的思想对预测施加安全余量（抬高负载预测、下压光伏预测）\cite{arrow1951}，并用全年仿真确定余量系数。
日内则利用储能对预测偏差进行实时纠偏：出现缺额时依次减少计划充电、增加储能放电，最后才紧急购电；
出现富余时依次减少计划放电、增加储能充电，储能调节到极限后才弃光；
日末储电量作为次日的初始状态，从而形成跨日连续运行的滚动调控策略。

\subsection{针对问题三的分析}
问题三在问题二的基础上增加了 0:00、6:00、12:00、18:00 的光伏预报，可据此滚动调整购电策略；
但偏离计划需付出代价（少购按 0.5 倍、多购按 1.5 倍结算），
因此关键不在于单纯追求预测精度，而在于权衡“预报更新带来的收益”与“调整成本、紧急购电风险”。
先对附件 3 的预报进行精度诊断：同一目标时刻的预报精度并不随发布时间改善，
且存在上午低估、午后高估的系统性偏差，故不宜直接用最新预报替换原计划，
而应结合历史同期信息对其标定，并利用当日已实现残差对后续时段进行修正。
在此基础上，以 0:00 制定的计划为基准，在 6:00、12:00、18:00 用更新后的预测重新优化剩余时段，
目标函数中同时计入正常购电费用、调整购电的违约/溢价费用与 5 倍紧急购电费用，
由模型自行判断“维持原计划”与“调整购电”哪一种更经济。
此外，若允许日末储电量完全自由，模型会在傍晚短视放空储能以降低当日费用，
故在调整模型中设置日末最低备用约束，并通过全年仿真确定其水平，
最终形成“预测标定—日内修正—滚动调整—末端备用”的调控策略。

\subsection{针对问题四的分析}
问题四进一步放开电价，使不同日期、不同时段的电价均可变化，
微网的购电决策除受负载、光伏与储能状态影响外，还取决于电价自身的时序变化，
因此需在问题二、问题三的模型框架内，将各项费用中的固定电价 $c_t$ 替换为当日第 $t$ 时段的实际电价 $c_{d,t}$ 重新求解。
先分析附件 4 的电价特征：全年平均水平与附件 1 相同（0.7662 元/kWh），
但取值范围明显扩大、日内价差更大，说明费用变化不能仅由均价解释，
还需进一步考察电价与负载、光伏在时间上的匹配关系。
其次，电价波动改变了储能的经济价值：日内价差越大，低充高放的能量时移收益越高；
但预测偏差的代价也随电价上行而放大，尤其在紧急购电按 5 倍结算时，
因此问题二与问题三两种策略对价格波动的适应性存在差异，需要比较其全年运行费用。


综上，问题四沿用问题二、问题三的基本模型结构，仅替换价格项并重新进行全年滚动优化，
通过对比固定电价与波动电价下的费用构成、紧急购电情况与储能周转量，
量化电价波动对微网购电与储能调度策略的影响。

\subsection{本文总体建模思路}
综合上述分析，四个问题共用同一套“微网功率平衡 $+$ 储能约束 $+$ 购电费用最小”的优化内核，
差异仅在于电价、负载与光伏信息的已知程度：问题一假定日内规律已知，求解日循环线性规划；
问题二在负载、光伏逐日变化的条件下引入日前预测、安全余量与日内实时纠偏；
问题三进一步利用整点光伏预报滚动调整购电计划，并设置日末最低备用约束；
问题四将固定电价替换为逐日波动的实际电价，重新求解并比较两类策略的经济性。
本文总体建模与求解的技术路线如图~\ref{fig:route} 所示。

\begin{figure}[!htbp]
\centering
\footnotesize
\begin{tikzpicture}[
  >=Stealth,
  data/.style={rectangle, rounded corners=1mm, draw=black!55, fill=blue!7, align=center,
               text width=28mm, minimum height=11mm, inner sep=2pt},
  core/.style={rectangle, draw=orange!80!black, fill=orange!12, thick, align=center,
               text width=124mm, minimum height=9mm, inner sep=2pt},
  prob/.style={rectangle, draw=black!65, fill=green!9, align=center,
               text width=28mm, minimum height=19mm, inner sep=2pt},
  outbox/.style={rectangle, rounded corners=1mm, draw=black!55, fill=gray!10, align=center,
               text width=124mm, minimum height=9mm, inner sep=2pt},
  line/.style={draw, -Latex, thick}
]
  \node[data] (d1) at (0,0)      {附件 1\\ 典型日电价、负载与光伏};
  \node[data] (d2) at (3.75,0)   {附件 2\\ 2025 年逐日负载与光伏};
  \node[data] (d3) at (7.5,0)    {附件 3\\ 整点光伏预报};
  \node[data] (d4) at (11.25,0)  {附件 4\\ 逐日波动电价};

  \node[core] (pre) at (5.625,-2.0) {数据预处理与规律分析：缺失值检查 $\cdot$ 负载星期周期 $\cdot$ 光伏持续性与季节性 $\cdot$ 预报精度诊断 $\cdot$ 电价特征统计};
  \node[core] (mod) at (5.625,-3.6) {共性优化内核：微网功率平衡 $+$ 储能容量/功率/效率约束 $+$ 购电费用最小（线性规划）};

  \node[prob] (q1) at (0,-5.7)     {问题一\\ 日内规律已知\\ 日循环线性规划\\ （首末储电量相等）};
  \node[prob] (q2) at (3.75,-5.7)  {问题二\\ w1 预测 $+$ 安全余量\\ 日前计划 $+$ 日内纠偏\\ 跨日滚动仿真};
  \node[prob] (q3) at (7.5,-5.7)   {问题三\\ 预报标定 $+$ 残差修正\\ 0/6/12/18 时滚动调整\\ 日末最低备用约束};
  \node[prob] (q4) at (11.25,-5.7) {问题四\\ 价格项 $c_t\!\rightarrow\! c_{d,t}$\\ 重现问题二、三策略\\ 费用构成对比};

  \node[outbox] (out) at (5.625,-7.9) {结果输出：指定时段（日期）购电量与费用 $\cdot$ 储能充放电量与储电量轨迹 $\cdot$ 紧急购电情况 $\cdot$ 全年费用对比与灵敏度分析};

  \draw[line] (d1.south) -- (d1 |- pre.north);
  \draw[line] (d2.south) -- (d2 |- pre.north);
  \draw[line] (d3.south) -- (d3 |- pre.north);
  \draw[line] (d4.south) -- (d4 |- pre.north);
  \draw[line] (pre.south) -- (mod.north);
  \draw[line] (mod.south) -- (5.625,-4.6);
  \draw[draw, thick] (0,-4.6) -- (11.25,-4.6);
  \draw[line] (0,-4.6) -- (q1.north);
  \draw[line] (3.75,-4.6) -- (q2.north);
  \draw[line] (7.5,-4.6) -- (q3.north);
  \draw[line] (11.25,-4.6) -- (q4.north);
  \draw[line] (q1.south) -- (q1 |- out.north);
  \draw[line] (q2.south) -- (q2 |- out.north);
  \draw[line] (q3.south) -- (q3 |- out.north);
  \draw[line] (q4.south) -- (q4 |- out.north);
\end{tikzpicture}
\caption{本文总体建模与求解技术路线}
\label{fig:route}
\end{figure}
